%% BLOQUE: datos
Se analizan 2011 muestras de agua con nueve mediciones fisicoquímicas (\texttt{ph},
dureza, sólidos disueltos, cloraminas, sulfato, conductividad, carbono orgánico,
trihalometanos y turbidez). El 59.7\,\% de las muestras es no potable. Un rasgo
central de este conjunto es que las medias de cada variable son casi idénticas en las
dos clases (por ejemplo, \texttt{ph} 7.07 frente a 7.11 y sulfato 333.7 frente a
332.5\,mg/L): ninguna variable separa las clases por sí sola. Lo que cambia es la
dispersión y los extremos. El 73\,\% de las muestras con sulfato menor que 260\,mg/L
y el 62\,\% de las que superan 420\,mg/L son potables, frente al 35\,\% entre 300 y
340\,mg/L. Con el \texttt{ph} ocurre algo parecido: la proporción de potables es mayor
cerca de la neutralidad (43--45\,\% entre 6 y 8) que en los extremos (30\,\% con
\texttt{ph} $\le 5$). Son relaciones en forma de U, no lineales, y eso condiciona qué
algoritmos funcionan mejor.

%% BLOQUE: KNN
Las cuatro configuraciones quedan dentro de una desviación estándar entre sí
(0.588--0.604, con desviaciones cercanas a 0.03), así que la elección de la distancia
Manhattan (+0.014 sobre el estándar) es una mejora pequeña y poco segura. Aumentar a
$k=15$ sube la exactitud de validación de 62.6\,\% a 65.7\,\% pero apenas el F1 macro:
promediar más vecinos suaviza la frontera hacia la clase mayoritaria. Ponderar por
distancia no cambia nada relevante. La selección de cinco atributos empeoró el F1 macro
a 0.536, por la razón que se discute en la síntesis. En prueba, KNN reconoce bien el
agua no potable (recall 0.787) pero sólo 83 de 203 muestras potables (recall 0.409).
La importancia por permutación reparte el efecto entre \texttt{ph}, sulfato y sólidos,
sin una variable dominante. Esto es coherente con el método: KNN usa las nueve
variables con el mismo peso, aunque la mayoría aporta poca información, y esas
dimensiones irrelevantes vuelven las distancias menos informativas. La diferencia entre
F1 de entrenamiento (0.747) y de prueba (0.597) muestra además que, con $k=5$, el
modelo depende mucho de los vecinos concretos.

%% BLOQUE: DT
El árbol sin restricciones (estándar) obtiene F1 macro 0.565 con una exactitud de
apenas 57.5\,\%, menor que la de predecir siempre ``no potable'': crece hasta
memorizar el entrenamiento. Limitar la profundidad a 5 sube la exactitud (62.5\,\%)
pero baja el F1 macro (0.535), porque un árbol tan corto termina prediciendo casi
siempre la clase mayoritaria. La mejor opción es exigir al menos 10 observaciones por
hoja (0.591, +0.025): regulariza sin impedir que el árbol aprenda las reglas en U. El
árbol final sigue siendo grande (profundidad 14, 107 hojas) y el F1 cae de 0.794 en
entrenamiento a 0.578 en prueba, lo que confirma su alta varianza. Sus primeras reglas
son interpretables y coinciden con lo observado en los datos: con sulfato
$\le 261$\,mg/L, sólidos mayores que 18\,347\,mg/L y \texttt{ph} $>6.23$ predice potable; con sulfato entre 261 y
387\,mg/L predice no potable, y por encima de 387\,mg/L vuelve a predecir potable
cuando \texttt{ph} $\le 7.91$. El árbol representa la forma en U combinando dos cortes
sobre la misma variable y una interacción con el \texttt{ph}.

%% BLOQUE: RF
Random Forest obtiene el mejor resultado de este dataset. El estándar (100 árboles)
alcanza 0.629 y pasar a 500 árboles sube el F1 macro a 0.644 y además baja la
desviación entre pliegues de 0.031 a 0.022: con más árboles el promedio es más estable,
que es exactamente lo que predice la teoría del bagging. Exigir 5 observaciones por
hoja empeora (0.615), porque limita los árboles profundos que el bosque necesita para
capturar las interacciones. Ponderar las clases no ayudó. En prueba logra exactitud de
70.2\,\% (contra 59.6\,\% de la referencia), F1 macro 0.659 y AUC 0.727. Es el modelo
más preciso al predecir potable (0.719: de cada 10 muestras que marca como potables,
7 lo son), pero sigue dejando pasar 116 de las 203 muestras potables (recall 0.429).
El F1 de entrenamiento es 1.000: cada árbol memoriza su muestra, y aun así el promedio
generaliza mejor que cualquier modelo individual. La importancia por permutación
(Figura~\ref{fig:perm-potabilidad}) identifica el sulfato (caída de 0.131 en F1) y el
\texttt{ph} (0.096) como las variables decisivas, seguidas de cloraminas, dureza y
sólidos.

%% BLOQUE: XGBoost
Las variantes de tasa de aprendizaje (0.05) y profundidad (3) quedan prácticamente
iguales al estándar (0.617--0.622). La única mejora es ponderar las clases (0.639,
+0.017), que obliga al modelo a prestar más atención a las muestras potables. En
prueba esto se ve como el segundo mayor recall de potables (0.478, 97 de 203), a costa
de más falsos positivos (90 frente a 34 de RF) y una exactitud de 61.0\,\%. XGBoost es
el modelo que más cae de validación (0.639) a prueba (0.590), y su F1 de entrenamiento
es 1.000: con la tasa 0.3 y profundidad 6 predeterminadas, los 100 árboles alcanzan a
memorizar el entrenamiento. En este dataset ruidoso, el boosting por gradiente termina
ajustando ruido, mientras que el promedio de RF lo atenúa.

%% BLOQUE: AdaBoost
AdaBoost estándar tiene el peor desempeño del estudio: F1 macro 0.471 con exactitud
59.7\,\%, igual a la proporción de la clase mayoritaria. Es decir, casi siempre
predice ``no potable''. Bajar la tasa a 0.5 lo empeora (0.450) y duplicar los
estimadores mejora poco (0.501). La razón está en sus aprendices: cada stump es un
único corte sobre una variable, y como las medias de ambas clases son casi iguales,
ningún corte supera por mucho el azar. En ese caso el peso $\alpha_t$ de cada stump es
pequeño y el ensamble no logra representar la interacción entre sulfato
y \texttt{ph}. La ponderación de clases es la variante con mayor efecto de todo el
dataset (+0.069, F1 0.540): al igualar el peso de ambas clases, los stumps dejan de
favorecer la mayoritaria. En prueba alcanza el recall más alto de potables (0.498),
pero con precisión 0.502 y AUC 0.615, la más baja. A diferencia de RF y XGBoost no
sobreajusta (F1 0.575 en entrenamiento y 0.582 en prueba): su problema es de sesgo
(subajuste), no de varianza. Casi toda su capacidad se apoya en el sulfato.

%% BLOQUE: sintesis
Por validación cruzada, el mejor modelo es Random Forest con 500 árboles (F1 macro
$0.644\pm0.022$), y también es el mejor en prueba en F1 macro, exactitud y AUC. Aun
así, el problema es difícil: las AUC van de 0.615 a 0.727 y ninguna familia detecta
más del 50\,\% de las muestras potables. Hay 37 muestras potables de prueba que los
cinco modelos clasifican como no potables. Con medias casi idénticas entre clases, la
información útil está en regiones extremas y en combinaciones de variables, y los
modelos que mejor la aprovechan son los ensambles de árboles profundos.

La selección de atributos por información mutua empeoró las cinco familias, con caídas
de 0.01 (AdaBoost) a 0.10 (RF) en F1 macro (Figura~\ref{fig:seleccion-potabilidad}). El motivo es
instructivo. La información mutua calculada en train asigna 0.026 al sulfato, 0.014 a
la dureza y prácticamente 0 al \texttt{ph} y a las cloraminas, así que el filtro
descarta el \texttt{ph}. Sin embargo, la importancia por permutación muestra que el
\texttt{ph} es la segunda variable más importante para RF. Su efecto aparece sobre todo
en combinación con el sulfato (como muestran las reglas del árbol), y un filtro que
evalúa cada variable por separado no lo detecta. En este dataset conviene conservar
las nueve variables y dejar que los árboles encuentren las interacciones.
